\documentclass[aspectratio=169, 14pt]{beamer}
\usepackage{beamer_preamble}

\begin{document}

% ==== Title ===============================================================
\titleslide{Lesson 4-3 \textperiodcentered\ Period 3}{Rational Expression Modeling + Assessment-Critical Simplify}

% ==== Do Now ==============================================================
\begin{frame}{Do Now --- Identity and Domain}{Algebra 2 \textperiodcentered\ Lesson 4-3 \textperiodcentered\ Period 3}
\phasebadges{2}{5}{bridge from P2}
\vspace{0.2cm}
\begin{projcard}{Do Now \textperiodcentered\ Identity check (4-3-savvas-q11)}{slidegold}
\textbf{Reason} Explain why
\[
\frac{4x^2-7}{4x^2-7}=1
\]
is a valid identity --- and state the values of \(x\) that must be restricted from the domain.
\end{projcard}

\vspace{0.15cm}
\begin{projcard}{Bridge prediction (30 sec, no wrong answers)}{slidegold}
Before Launch: what do you think the SA/Volume ratio simplifies to for a cube with side \(x\)?\\
Write a guess --- we compare in Launch.
\end{projcard}
\end{frame}

% ==== Launch ==============================================================
\begin{frame}{Launch --- Example 2 + Example 6 (back-to-back)}{Algebra 2 \textperiodcentered\ Lesson 4-3 \textperiodcentered\ Period 3}
\phasebadges{2}{12}{teacher models Ex 2 + Ex 6}
\vspace{0.15cm}
\begin{projcard}{DOMAIN RULES (mark these --- assessment reference)}{slidegreen}
\begin{enumerate}
  \setlength{\itemsep}{0pt}
  \item Domain = all reals EXCEPT zeros of \textbf{every} denominator in every step.
  \item Write the simplified form AND the restriction list together.
  \item Canceled factors still exclude their zeros.
\end{enumerate}
\end{projcard}
\vspace{-0.1cm}
\begin{projcard}{MODELING RULES (mark these too)}{slidegreen}
\begin{enumerate}
  \setlength{\itemsep}{0pt}
  \item Identify the ratio in words first (what / what).
  \item Translate each quantity to a rational expression, with units.
  \item Simplify, restrict domain, then interpret the simplified ratio in context.
\end{enumerate}
\end{projcard}

\vspace{0.05cm}
\small\textcolor{slidegray}{Ex 2: simplify \((4-x^2)/(x^2+3x-10)\) \textbullet\ Ex 6: SA/Volume for prism vs.\ cylinder. Both land on \(3/x\).}
\end{frame}

% ==== Explore =============================================================
\begin{frame}{Explore --- TPS: 6 items in order (35 min)}{Algebra 2 \textperiodcentered\ Lesson 4-3 \textperiodcentered\ Period 3}
\phasebadges{2}{35}{student work}
\small\textcolor{slidegray}{Abstract forms first, then modeling, then assessment rehearsal. Never cut \#19.}

\vspace{0.1cm}
\begin{projcard}{Try It 6 \textperiodcentered\ SA/V comparison (two prisms)}{slideblue}
Which container has the smaller surface area-to-volume ratio?
\end{projcard}
\vspace{-0.12cm}
\begin{projcard}{Practice \#14 \textperiodcentered\ Domain restriction reasoning}{slideblue}
Why are restrictions needed to show \(\dfrac{-6x^2+21x}{3x}\) and \(-2x+7\) are equivalent?
\end{projcard}
\vspace{-0.12cm}
\begin{projcard}{Practice \#22 \textperiodcentered\ Simplify + identify domain}{slideblue}
Simplify \(\dfrac{y^2-5y-24}{y^2+3y}\). State the domain.
\end{projcard}
\vspace{-0.12cm}
\begin{projcard}{Practice \#34 \textperiodcentered\ Volume ratio (two prisms)}{slideblue}
Find and simplify the ratio of Volume A to Volume B.
\end{projcard}
\vspace{-0.12cm}
\begin{projcard}{Practice \#37 \textperiodcentered\ Parallelogram base from area}{slideblue}
Area \(= \dfrac{3x+12}{10x+25}\), height shown. Find the base length.
\end{projcard}
\vspace{-0.12cm}
\begin{projcard}{Practice \#19 \textperiodcentered\ \(\bigstar\) ASSESSMENT REHEARSAL --- Topic 4 LEHS Q\#2}{slideaccent}
Denominator \(= x^3+3x^2-10x\). What values must be restricted from the domain?
\end{projcard}
\end{frame}

% ==== Share / Summary =====================================================
\begin{frame}{Share / Summary --- Assessment-Critical Reveal}{Algebra 2 \textperiodcentered\ Lesson 4-3 \textperiodcentered\ Period 3}
\phasebadges{2}{5}{DOK-2 reveal + Topic 4 assessment preview}
\vspace{0.15cm}
\begin{projcard}{Sentence frame for \#19 (pair presents, class signals)}{slideblue}
``The denominator \(x^3+3x^2-10x\) factors as \underline{\hspace{2.5cm}}.\\
The zeros are \underline{\hspace{1cm}}, \underline{\hspace{1cm}}, \underline{\hspace{1cm}}, so \(x \ne\) \underline{\hspace{2.5cm}}.\\
A cubic denominator may have up to \underline{\hspace{0.6cm}} restrictions.''
\end{projcard}

\vspace{0.15cm}
\begin{projcard}{Topic 4 8-Q Assessment --- Q\#2 is this exact content}{slideaccent}
\textbf{Practice \#19 directly rehearses Topic 4 LEHS Q\#2.}\\
Assessment cue: cubic denominator \(\rightarrow\) factor completely \(\rightarrow\) list ALL zeros.\\
Do NOT stop at one restriction.
\end{projcard}
\end{frame}

% ==== Exit Ticket =========================================================
\begin{frame}{Exit Ticket --- Pre-Assessment Confidence Check}{Algebra 2 \textperiodcentered\ Lesson 4-3 \textperiodcentered\ Period 3}
\phasebadges{1-2}{3}{not graded}
\vspace{0.2cm}
\begin{projcard}{Complete on your exit ticket}{slidegold}
\begin{enumerate}
  \setlength{\itemsep}{0.25cm}
  \item A rational model needs factoring because \underline{\hspace{4.5cm}}.
  \item My simplified ratio is \underline{\hspace{2.5cm}}, valid when \(x \ne\) \underline{\hspace{1.5cm}}.
  \item In context this means \underline{\hspace{4.5cm}}.
  \item One biggest thing I learned today: \underline{\hspace{4cm}}.
\end{enumerate}
\end{projcard}
\end{frame}

\end{document}
